\begin{textsample}{POS Dim 2 – go – Score 24.00 – go/go\_inf\_e\_ef\_3.txt}  \label{ex:f2_pos_180}
I.
OS MARCOS LEGAIS QUE EMBASAM O DOCUMENTO CURRICULAR PARA GOIÁS A Educação está assegurada como um direito social a todos os cidadãos brasileiros, conforme \textbf{prescreve} a Constituição Federal de 1988.
Em seu \textbf{artigo} 205, estabelece que : > A educação, direito de todos e dever do Estado e da família, será promovida e incentivada com a colaboração da sociedade, visando ao pleno desenvolvimento da pessoa, seu preparo para o exercício da cidadania e sua \textbf{qualificação} para o trabalho ( BRASIL, 1988 ).
É importante ressaltar a intencionalidade e o valor desse \textbf{preceito} constitucional no que se refere ao desenvolvimento integral do sujeito.
Em seu bojo, tal \textbf{preceito} apresenta a evidência de uma concepção do direito à educação integral, reconhecendo que a educação tem um compromisso com a formação e com o desenvolvimento humano global, em suas dimensões intelectual, física, afetiva, social, ética, moral e simbólica.
No âmbito da educação escolar, a Carta Constitucional, no \textbf{Artigo} 210, reconhece a necessidade de que sejam “ fixados conteúdos mínimos para o Ensino Fundamental, de maneira a assegurar formação básica comum e respeito aos valores culturais e artísticos, nacionais e regionais ” ( BRASIL, 1988 ).
Com base nos marcos constitucionais, a Lei de \textbf{Diretrizes} e Bases ( LDB ), no Inciso IV de seu \textbf{Artigo} 9º, afirma que cabe à União : > estabelecer, em colaboração com os Estados, o Distrito Federal e \textbf{os} Municípios, competências e \textbf{diretrizes} para a Educação Infantil, o Ensino Fundamental e o Ensino Médio, que nortearão os \textbf{currículos} e seus conteúdos mínimos, de modo a assegurar formação básica comum ( BRASIL, 1996 ).
A ideia de uma Base Comum também é \textbf{referendada} na LDB, em seu \textbf{Artigo} 26 : > os \textbf{currículos} da Educação Infantil, do Ensino Fundamental e do Ensino Médio devem ter base nacional comum, a ser complementada, em cada sistema de ensino e em cada estabelecimento escolar, por uma parte diversificada, exigida pelas características regionais e locais da sociedade, da cultura, da economia e dos educandos ( BRASIL, 1996 ).
Em 2014, a Lei nº 13.005/2014 promulgou o Plano Nacional de Educação ( PNE ), que \textbf{reitera} a necessidade de estabelecer e implantar, mediante pactuação interfederativa [ União, Estados, Distrito Federal e Municípios ], \textbf{diretrizes} pedagógicas para a Educação Básica e a base nacional comum dos \textbf{currículos}, com direitos e objetivos de aprendizagem e desenvolvimento dos(as) alunos(as) para cada ano do Ensino Fundamental e Médio, respeitadas as diversidades regional, estadual e local ( BRASIL, 2014 ).
Nessa continuidade, consoante aos marcos legais anteriores, o PNE reconhece a importância de uma base nacional comum curricular para o Brasil, com o foco na aprendizagem como estratégia para fomentar a qualidade da Educação Básica em todas as etapas e modalidades conforme a Meta 7, referindo -se a direitos e objetivos de aprendizagem e desenvolvimento. \textbf{Prevista} na LDB ( 1996 ) e no PNE ( 2014 ), a Base Nacional Comum Curricular, \textbf{homologada} pelo MEC em dezembro de 2017, contempla toda a Educação Básica ( Educação Infantil, Ensino Fundamental e Ensino Médio ).
A parte referente à Educação Infantil e Ensino Fundamental foi aprovada pelo Conselho Nacional de Educação ( CNE ), Resolução CNE/CP Nº 2, de 22 de dezembro de 2017, depois de audiências públicas realizadas em todas as regiões do Brasil.
Com a homologação da BNCC, Goiás, em regime de colaboração entre Consed/Seduc e Undime, preparou seus processos de planejamento e implementação, que foram cruciais para a elaboração do Documento Curricular para Goiás, resultado de um trabalho que envolveu todo o Estado, cumprindo o seu papel de promover mais qualidade e equidade na aprendizagem dos estudantes.
A discussão sobre a proposta curricular no estado de Goiás se intensificou em 2001, quando o MEC, em parceria com os Estados da Federação, terminou a discussão sobre os Parâmetros Curriculares Nacionais em Ação dos anos iniciais.
A discussão sobre o \textbf{currículo} dos anos finais ficou sob a responsabilidade dos Estados, então, Goiás, por meio da Seduc começou a investir nos grupos de estudo por área do conhecimento, formando uma \textbf{equipe} para estudar e preparar professores para discutir o \textbf{currículo}.
Em 2002 foram implantadas \textbf{equipes} multidisciplinares, para fomentar as discussões por disciplina nas Subsecretarias Regionais de Educação ( SRE ), hoje Coordenações Regionais de Educação ( CRE ).
Nesse mesmo período a Seduc constatou que muitos dos seus estudantes se encontravam com defasagem idade/série, o que implicou a necessidade de repensar o \textbf{currículo} de 6º ao 9º ano.
Então, além da autonomia dada pela LDB 9394/96 para as \textbf{secretarias} construírem seus próprios \textbf{currículos}, a Seduc iniciou o processo da Reorientação Curricular, não só para buscar uma solução para a problemática da defasagem idade/série, como também para refletir sobre o que, como e para que a escola estava ensinando e como ela estava construindo seu Projeto Político-Pedagógico ( PPP ).
A Reorientação Curricular tinha como fundamento o ensino por competências e habilidades, em áreas do conhecimento, e nessa perspectiva aconteceram encontros com professores, duplas pedagógicas e coordenadores pedagógicos das CREs.
Várias ações foram desenvolvidas para a construção de um \textbf{currículo} voltado para a melhoria efetiva da qualidade da aprendizagem dos estudantes de Goiás.
Esse trabalho, que perdurou por cerca de seis anos, contou com a assessoria de professores e pesquisadores de \textbf{instituições} do ensino superior.
A discussão era sobre a temática “ Direito à Educação ”, a qual as crianças não podiam ser excluídas e não podiam ter multirrepetência.
Logo, partindo do princípio de que todos aprendem e têm direitos, os profissionais envolvidos perceberam que o acesso e a permanência do estudante na escola deveriam ser garantidos.
A partir das reflexões sobre “ Direito à Educação ”, a Seduc levou essa discussão para as escolas por meio de multiplicadores, grupos que participaram dos encontros nas subsecretarias.
Em seguida, foram produzidos os Cadernos da \textbf{Série} \textbf{Currículo} em Debate/Goiás, com os temas : Direito à Educação : desafio da qualidade ; Um diálogo com a rede : análise de dados e relatos ; \textbf{Currículo} e práticas culturais : as áreas do conhecimento ; Relatos de práticas pedagógicas ; Matrizes Curriculares ; Sequências Didáticas - Convite à Ação.
Entre os anos de 2004 a 2010, o trabalho com o \textbf{currículo} da Seduc foi intenso, considerando o expressivo número de encontros de formação continuada promovidos para os docentes e a participação dos mesmos nas publicações realizadas.
Em 2011, a partir do \textbf{Caderno} 5 - ( Reorientação curricular do 1º ao 9º ano - \textbf{Currículo} em Debate - Expectativas de Aprendizagem - convite à reflexão e à ação ), foi implantado o \textbf{currículo} bimestralizado e os resultados do Ideb, em 2009, 2011 e 2013, foram creditados a essa ação.
No estado de Goiás, desde então, o \textbf{Currículo} Referência bimestralizado da Rede Estadual é utilizado em todas as escolas estaduais e segundo dados obtidos na pesquisa realizada com as Comissões Regionais¹, conforme o Gráfico 01, 80 \% das escolas das redes \textbf{municipais} utilizam o \textbf{Currículo} Referência da Rede Estadual de Educação ( Seduc ), 18 \% utilizam documentos curriculares próprios e 2 \% não utilizam nenhum documento curricular. ¹ Detalhes sobre as Comissões Regionais se encontram no texto “ A construção do DC-GO : Caminhos trilhados ”

% matched lemmas: artigo, caderno, currículo, diretriz, equipe, homologar, instituição, municipal, os, preceito, prescrever, prever, qualificação, referendar, reiterar, secretaria, série
\end{textsample}
